\chapter{{\fontsize{14pt}{21pt}\selectfont\MakeUppercase{Данные проекта}}}
\label{ch:data}

\section{Какие данные используются в проекте}

В~проекте используются следующие группы данных.

\textbf{Реляционные данные (PostgreSQL):} пользователи (учётные записи, роли \texttt{owner}/\texttt{admin}/\texttt{member}), бизнес-объекты (название, категория, адрес, телефон, описание, расписание работы), интеграции (платформа, статус, зашифрованные токены доступа, внешний идентификатор), подписки (план, статус, срок действия), журнал аудита (действия, ресурс, детали), токены обновления JWT, статусы синхронизации по платформам (для обнаружения расхождений).

\textbf{Документные и слабоструктурированные данные (MongoDB):} диалоги с~оркестратором (история сообщений, вызовы инструментов, результаты); задачи агентов (тип, статус, платформа, вход/выход, ошибки); отзывы с~платформ (разная схема по платформам: автор, рейтинг, текст, ответ); публикации (контент, медиа, результаты по платформам, статус, дата публикации).

\section{Откуда берутся данные}

Данные поступают из следующих источников.

\begin{itemize}
  \item \textbf{Ввод пользователя:} регистрация, профиль бизнеса, расписание, настройки; подключение каналов (OAuth для VK, токен бота для Telegram, сессия для Яндекс.Бизнес).
  \item \textbf{API внешних платформ:} VK API и Telegram Bot API~--- получение и публикация контента, комментарии, отзывы (где доступно).
  \item \textbf{RPA (браузерная автоматизация):} Яндекс.Бизнес~--- чтение и обновление карточки организации, отзывы и ответы на них через веб-интерфейс.
  \item \textbf{Внутренние процессы:} оркестратор генерирует сообщения и вызовы инструментов; результаты исполнения и логи сохраняются в~БД.
\end{itemize}

Данные открытых API используются в~соответствии с~условиями использования платформ; авторские права на контент остаются у владельцев. Персональные данные пользователей (учётные записи, привязка к~бизнесу) обрабатываются в~рамках проекта.

\section{Работа с персональными данными}

Обработка персональных данных осуществляется в~соответствии с~требованиями законодательства Российской Федерации (152-ФЗ). Хранение учётных данных (хеш пароля, токены доступа к~платформам) выполняется с~применением криптографической защиты: пароли~--- bcrypt; токены платформ~--- шифрование AES-256-GCM при хранении; ключи шифрования не хранятся в~БД. Реализованы разграничение доступа по ролям (RBAC), аудит значимых действий (журнал аудита), JWT с~ограниченным сроком жизни для доступа к~API. В~логах не выводятся чувствительные данные; идентификация запросов~--- по корреляционному идентификатору. Требования к~информационной безопасности и совместимости зафиксированы в~документации проекта и техническом задании.
